\begin{table}[htbp]
\centering
\caption{Class size and test scores: original synthetic data}
\label{tab:lab04-class-size}
\begingroup
\small
\setlength{\tabcolsep}{6pt}
\renewcommand{\arraystretch}{1.12}
\begin{adjustbox}{max width=\linewidth}
\begin{tabular}{@{}lcc@{}}
\toprule
Dependent variable: & test\_score & test\_score \\
 & (1) & (2) \\
 & HC3 robust & Classical \\
\midrule
Students per class & $-1.943$ & $-1.943$ \\
 & $(0.145)$ & $(0.142)$ \\[0.35em]
Constant & $705.723$ & $705.723$ \\
 & $(3.200)$ & $(3.424)$ \\[0.35em]
\midrule
Observations & 240 & 240 \\
$R^{2}$ & 0.442 & 0.442 \\
$\text{Adjusted }R^{2}$ & 0.440 & 0.440 \\
\bottomrule
\end{tabular}
\end{adjustbox}
\par\smallskip
\begin{minipage}{\linewidth}
\footnotesize
\raggedright
\textit{Notes:} Standard errors are in parentheses.\par
(1) OLS: HC3 heteroskedasticity-robust standard errors; finite-sample covariance factor = 1; Student t inference (df = 238); 95\% confidence level (alpha = 0.05); estimation sample 240 of 240 observations; model F(1, 238) = 180.591, p = 5.1645e-31.\par
(2) OLS: conventional standard errors; Student t inference (df = 238); 95\% confidence level (alpha = 0.05); estimation sample 240 of 240 observations; model F(1, 238) = 188.462, p = 5.56096e-32.\par
Identical 240-school sample in both columns; only covariance differs. School outcomes are simulated, not empirical policy evidence.\par
\end{minipage}
\endgroup
\end{table}